% Avsnitt 1.8, ur lectures/L01/appendix/b_exercises.md.

\section{Övningar}
\label{c1:sec:exercises}

\begin{aside}
\textbf{Så kontrollerar du ditt arbete.} Övning~\ref{c1:ex:sop} är en kort uppvärmning med papper
och penna på den notation boken använder rakt igenom; den kräver varken VHDL, GHDL eller kort, och
anger vad ett bra svar täcker.

Från Övning~\ref{c1:ex:gates} och framåt har varje övning som ber dig skriva en VHDL-modul en
utdelad självkontrollerande testbänk under \file{lectures/L01/exercises/}. Skriv din modul i dess
katalog \file{exercises/<module>/}, med det entitetsnamn och den \textbf{portordning} övningen
anger, och kör den sedan med GHDL \textemdash{} se \secref{c2:sec:testbenches} för de tre
kommandona.

Inget FPGA-kort behövs till någon övning. Stegen med Quartus-syntes och kortprogrammering
demonstreras under passet; ditt jobb efteråt är att få VHDL-koden rätt, och testbänken är hur du
bekräftar det.
\end{aside}

% ------------------------------------------------------------------------------------------
\exerciseset{Sanningstabeller och grindar}

\exercise{En funktion som beror på mindre än den ser ut att göra}{Krets}
\label{c1:ex:sop}
En funktion av två ingångar \code{A}, \code{B} och en utgång \code{X} har sanningstabellen nedan:

\begin{narrowtable}{ccc}
  \thead{A} & \thead{B} & \thead{X} \\
  \midrule
  0 & 0 & 1 \\
  0 & 1 & 0 \\
  1 & 0 & 1 \\
  1 & 1 & 0 \\
\end{narrowtable}

\xpart{a} Härled summa-av-produkter-ekvationen för \code{X} direkt ur sanningstabellen
(\secref{c1:sec:boolean}). Två rader har $X = 1$, så du får två AND-termer.

\xpart{b} Förenkla den algebraiskt. Bryt ut det som de två termerna delar, och tillämpa
$A + A' = 1$ och $1 \cdot B' = B'$ från \secref{c1:sec:boolean}. Du ska hamna i en enda literal.

\xpart{c} Ditt svar på \textbf{b)} säger något om \code{A}. Vad? Vilken enda grind ur tabellen i
\secref{c1:sec:gates} beräknar funktionen, och vad händer med ingången \code{A} när du ritar den?

\xpart{d} Bygg den enkelgrindsversionen i CircuitVerse och bekräfta att den återger varje rad i
sanningstabellen ovan. Lämna ingången \code{A} okopplad, eller utelämna den helt enkelt ur
ritningen; \textbf{b)} är motiveringen till att göra så.

% ------------------------------------------------------------------------------------------
\exerciseset{Entiteter och arkitekturer}

\exercise{\code{and_gate} och \code{nand_gate}}{VHDL}
\label{c1:ex:gates}
Skriv två små grindar som kompletta VHDL-moduler, \code{and_gate} och \code{nand_gate}.

Ingen av grindarna är poängen, och du ska inte förvänta dig att någon av dem är svår. Poängen är
att det är här en entitet och en arkitektur slutar vara något du läser om i \secref{c1:sec:vhdl}
och blir något du skriver, och att de två halvorna är värda att skriva i den ordningen.

Båda entiteterna har samma portar:

\begin{narrowtable}{lll}
  \thead{Port} & \thead{Riktning} & \thead{Typ} \\
  \midrule
  \code{a} & in & \code{std_logic} \\
  \code{b} & in & \code{std_logic} \\
  \code{x} & out & \code{std_logic} \\
\end{narrowtable}

\xpart{a} Skriv \code{entity}-deklarationen för \code{and_gate}, och ingenting annat. I VHDL är en
moduls gränssnitt något du kan skriva ner och resonera om innan någon implementation alls finns,
vilket är den vana den här deluppgiften finns till för att bygga. Identifiera vilka portar som
krävs, riktningen på var och en, och typen på var och en.

\xpart{b} Lägg till en \code{architecture} som implementerar AND-grinden med en enda konkurrent
signaltilldelning.

\xpart{c} Skriv nu \code{nand_gate}: samma gränssnitt med inverterad utgång, med VHDL:s
\code{nand}-operator direkt. Skriv \textbf{inte} \code{not (a and b)}.

De två är likvärdiga som boolesk algebra, och på en FPGA gör de ingen skillnad alls, eftersom båda
hamnar i samma lookup-tabell. Det spelar roll på den sortens hårdvara som Övning~\ref{c1:ex:nand}
nedan handlar om, där NAND är den grind du fysiskt har och allt annat måste byggas av den. Att
skriva vad du menar, och låta verktygen avgöra vad det kostar, är vanan; att sträcka sig förbi en
operator språket redan ger dig är det som ska undvikas.

\modulefig[0.52]{lectures/L01/appendix/images/and_gate.png}
\modulefig[0.52]{lectures/L01/appendix/images/nand_gate.png}

\begin{selfcheck}
\textbf{Självkontroll.} Döp dina entiteter till \code{and_gate} och \code{nand_gate}, var och en
med portarna \code{a}, \code{b} (in) och \code{x} (out), deklarerade i den ordningen; deras
testbänkar ligger i \file{exercises/and_gate/} och \file{exercises/nand_gate/}, och var och en
kontrollerar alla fyra raderna i motsvarande sanningstabell.
\end{selfcheck}

\exercise{\code{xor3}: udda paritet med tre ingångar}{VHDL}
\label{c1:ex:xor3}
En XOR-grind har tre ingångar, \code{a}, \code{b}, \code{c}, och en utgång, \code{x}. Utgången är
\code{1} när ett udda antal av ingångarna är \code{1} (XOR med tre ingångar, eller
paritetsfunktionen).

\xpart{a} Skriv sanningstabellen för \code{x} över alla åtta kombinationer av \code{a}, \code{b},
\code{c}, och bekräfta att den stämmer med beskrivningen ”ett udda antal \code{1}:or”.

\xpart{b} Implementera grinden i VHDL, som en egen modul med en entitet vars portar är tre
\code{std_logic}-ingångar och en \code{std_logic}-utgång, och en arkitektur som driver \code{x}
med en enda konkurrent tilldelning. \Secref{c1:sec:vhdl} visar mönstret entitet/arkitektur och
noterar att VHDL:s \code{xor}-operator arbetar direkt på \code{std_logic};
\secref{c1:sec:module} visar en komplett modul från början till slut.

\xpart{c} Bygg samma grind i CircuitVerse och bekräfta att den återger varje rad i din
sanningstabell.

\note[Tips]\code{xor} kedjar från vänster till höger, så hela arkitekturkroppen är en enda rad,
\code{x <= a xor b xor c;}, utan behov av vare sig intern \code{signal} eller mellanliggande
grindar.

\modulefig[0.52]{lectures/L01/appendix/images/xor3.png}

\begin{selfcheck}
\textbf{Självkontroll.} Döp din VHDL-entitet till \code{xor3}, med portarna \code{a}, \code{b},
\code{c} (in) och \code{x} (out), deklarerade i den ordningen; dess testbänk ligger i
\file{exercises/xor3/}, och \secref{c1:sec:module} listar de tre kommandon som kör den.
\end{selfcheck}

\exercise{\code{majority3}: en trippelredundant omröstning}{VHDL}
\label{c1:ex:majority3}
En \textbf{majoritetsgrind} har tre ingångar, \code{a}, \code{b}, \code{c}, och en utgång,
\code{x}. Utgången är \code{1} när \emph{minst två} av de tre ingångarna är \code{1}.

Det här är kretsen bakom en trippelredundant omröstning: tre sensorer rapporterar samma mätvärde,
och systemet agerar på det svar två av dem är överens om, så att en enskild trasig sensor inte kan
bestämma något på egen hand.

\xpart{a} Skriv sanningstabellen för \code{x} över alla åtta kombinationer av \code{a}, \code{b},
\code{c}.

\xpart{b} Härled summa-av-produkter-ekvationen för \code{x} direkt ur sanningstabellen
(\secref{c1:sec:boolean}):
\begin{itemize}
  \item Fyra rader har $x = 1$, så du får fyra AND-termer.
  \item Varje term innehåller alla tre ingångarna, primmade där ingången är \code{0} på sin rad.
\end{itemize}

\xpart{c} Implementera grinden i VHDL, som en egen modul:
\begin{itemize}
  \item En entitet med tre \code{std_logic}-ingångar och en \code{std_logic}-utgång.
  \item En arkitektur som driver \code{x} med en enda konkurrent tilldelning, och som översätter
    dina fyra AND-termer rakt av (\secref{c1:sec:vhdl}).
\end{itemize}

\xpart{d} Bygg samma nät i CircuitVerse och bekräfta att det återger varje rad i din
sanningstabell.

\note[Tips]Skriv ut ekvationen på papper innan du öppnar vare sig editorn eller CircuitVerse. Fyra
AND-termer med tre ingångar plus en OR-grind med fyra ingångar är mycket ledningsdragning att
rätta i efterhand.

\note[Blick framåt]Räkna grindarna din ekvation kräver, och övertyga dig sedan genom inspektion om
att $x = ab + ac + bc$ beräknar samma funktion med tre AND-grindar med två ingångar i stället för
fyra med tre. Summa-av-produkter gav dig ett \emph{korrekt} nät, inte det \emph{minsta}; att hitta
det mindre systematiskt, snarare än genom att stirra på det, är vad \chapref{c2:ch} är till för.

\modulefig[0.52]{lectures/L01/appendix/images/majority3.png}

\begin{selfcheck}
\textbf{Självkontroll.} Döp din VHDL-entitet till \code{majority3}, med portarna \code{a},
\code{b}, \code{c} (in) och \code{x} (out), deklarerade i den ordningen; dess testbänk ligger i
\file{exercises/majority3/} och kontrollerar alla åtta raderna. Vilken som helst av ekvationerna
går igenom \textemdash{} testbänken kontrollerar funktionen, inte vilket nät du valde.
\end{selfcheck}

\exercise{\code{half_adder}: den första modulen med två utgångar}{VHDL}
\label{c1:ex:half_adder}
Skriv en \textbf{halvadderare}: \code{sum} och \code{carry} ur två enskilda bitar. Du kan kretsen;
skälet till att den står här är att den är bokens första modul med \textbf{två utgångar}, och det
är en VHDL-fråga snarare än en logikfråga.

\xpart{a} Implementera den, med två konkurrenta tilldelningar och utan intern \code{signal}.

Två utgångar betyder inte två entiteter, två arkitekturer eller en process. Varje utgångsport
drivs av sin egen konkurrenta tilldelning, och de två är inte steg som sker i en ordning: de
beskriver två separata ledningsstycken, båda aktiva hela tiden. Om du kommer på dig själv med att
sträcka dig efter en process för att ”returnera” två värden är den instinkten
mjukvaruinstinkten, och \secref{c1:sec:vhdl}:s ”en konkurrent tilldelning beskriver en ledning”
är rättelsen.

\xpart{b} Bygg den i CircuitVerse och bekräfta båda utgångarna, så att du har sett formen med två
utgångar som en ritning innan du möter den som en entitet med flera \code{out}-portar i
\chapref{c2:ch}.

\note[Blick framåt]Det här är cellen som räknaren i \chapref{c6:ch} är byggd av. Adderaren med
fyra bitar där är fyra sådana kedjade carry-till-carry, och räknaren är den adderaren med sin
utgång återkopplad till sin ingång, vilket är därför överslaget inte behöver någon logik alls.

\modulefig[0.56]{lectures/L01/appendix/images/half_adder.png}

\begin{selfcheck}
\textbf{Självkontroll.} Döp din entitet till \code{half_adder}, med portarna \code{a}, \code{b}
(in) och \code{sum}, \code{carry} (out), deklarerade i den ordningen; dess testbänk ligger i
\file{exercises/half_adder/} och kontrollerar \code{sum} och \code{carry} var för sig, så att en
konstruktion som har den ena rätt och den andra fel får veta vilken som är vilken.
\end{selfcheck}

% ------------------------------------------------------------------------------------------
\exerciseset{Universella grindar}

\exercise{\code{or_from_nand}}{VHDL}
\label{c1:ex:nand}
\Secref{c1:sec:gates} slår fast att varje boolesk funktion kan byggas med enbart
\code{NAND}-grindar. Den här övningen får dig att bevisa det för dig själv för de tre operatorer
du har använt.

\xpart{a} Härled \code{OR}-fallet på papper innan du bygger något, eftersom De Morgan
(\secref{c1:sec:boolean}) ger dig det direkt. Utgå från $(A'B')'$ och tillämpa lagen; du ska komma
fram till $A + B$. Läs nu tillbaka det uttrycket som en krets: vad är $A'B'$ med en inversion på
sin utgång, och hur många \code{NAND}-grindar är alltihop?

\xpart{b} Bygg i CircuitVerse var och en av följande av enbart \code{NAND}-grindar, och stäm av
var och en mot dess sanningstabell i \secref{c1:sec:gates} efter hand:
\begin{itemize}
  \item \code{NOT a}: en NAND-grind räcker, om du matar samma signal till båda dess ingångar.
  \item \code{a AND b}: en NAND-grind, följd av den \code{NOT} du just byggde.
  \item \code{a OR b}: nätet du härledde i \textbf{a)}. Stäm av ditt grindantal mot vad du
    förutsade.
\end{itemize}

\xpart{c} Skriv \code{OR}-versionen i VHDL, som en egen modul:
\begin{itemize}
  \item En entitet med två \code{std_logic}-ingångar och en \code{std_logic}-utgång.
  \item En arkitektur som driver \code{x} med en enda konkurrent tilldelning.
  \item Använd enbart VHDL:s \code{nand}-operator: ingen \code{or}, ingen \code{and}, ingen
    \code{not} någonstans i uttrycket.
\end{itemize}

\xpart{d} Förklara varför det här spelar roll i praktiken. Tänk på att en chiptillverkningsprocess
fysiskt måste implementera varje grindtyp den erbjuder, och att \code{NAND} är den billigaste
grinden att bygga i CMOS.

\note[Tips]VHDL:s \code{nand}-operator kedjar inte så som \code{and} och \code{or} gör:
\code{a nand b nand c} är inte giltigt. Parentessätt varje NAND explicit, vilket ändå är vad du
vill här, eftersom varje parentespar är en fysisk grind i din CircuitVerse-ritning.

\modulefig[0.52]{lectures/L01/appendix/images/or_from_nand.png}

\begin{selfcheck}
\textbf{Självkontroll.} Döp din VHDL-entitet till \code{or_from_nand}, med portarna \code{a},
\code{b} (in) och \code{x} (out), deklarerade i den ordningen; dess testbänk ligger i
\file{exercises/or_from_nand/} och kontrollerar den mot alla fyra raderna i
\code{OR}-sanningstabellen.
\end{selfcheck}

% ------------------------------------------------------------------------------------------
\exerciseset{Att sätta ihop det hela}

\exercise{Bromsassistenten i VHDL}{VHDL}
\label{c1:ex:adas}
Skriv bromsassistenten från \secref{c1:sec:adas} i VHDL.

\Secref{c1:sec:adas} lät dig härleda dess ekvationer, räkna dess grindar och skissa dess nät
\emph{före} passet, och passet byggde den sedan live. Det är här du skriver den själv, och där en
testbänk avgör om din läsning av kravet stämmer med den i sanningstabellen.

Entiteten måste ha:

\begin{booktable}{@{}p{8em}p{3.4em}p{5.6em}L@{}}
  \thead{Port} & \thead{Riktn.} & \thead{Typ} & \thead{Beskrivning} \\
  \midrule
  \code{driver_brake} & in & \code{std_logic} & Föraren står på bromspedalen. \\
  \code{sensor} & in & \code{std_logic} & Avståndssensorn ser ett hinder. \\
  \code{radar} & in & \code{std_logic} & Radarn ser ett hinder. \\
  \code{error} & in & \code{std_logic} & Assistanssystemet rapporterar ett fel. \\
  \code{engine_brake} & out & \code{std_logic} & Lägg an bromsarna. \\
\end{booktable}

\xpart{a} Implementera den utifrån \textbf{dina egna} ekvationer snarare än utifrån det som stod
på skärmen. Om de två är oense är den oenigheten det mest användbara den här övningen kan ge dig,
och testbänken talar om vilken av dem sanningstabellen håller med.

Använd två konkurrenta tilldelningar och en intern \code{signal}:
\begin{itemize}
  \item Signalen håller assistanssystemets eget beslut, innan förarens pedal vägs in.
  \item Det här är bokens första övning med en intern signal, och det är anledningen till att
    nyckelordet finns: ett namn på ett mellanresultat, så att varje rad fortfarande läser sig som
    den sats i kravet den kom ifrån.
\end{itemize}

\xpart{b} Gör nu sönder den avsiktligt, på det specifika sätt \secref{c1:sec:adas} varnar för.
Mata \code{driver_brake} \emph{genom} fellogiken i stället för att låta den gå förbi:

\begin{vhdlcode}
engine_brake <= (driver_brake or sensor or radar) and (not error);
\end{vhdlcode}

Det är en grind enklare, och det ser rimligt ut. Kör testbänken mot den och svara:
\begin{itemize}
  \item Vilken ingångskombination faller den på först, och vad händer fysiskt i bilen i det
    ögonblicket?
  \item Hur många av de sexton raderna får den fel, och vad har de alla gemensamt?
  \item De rader som fallerar är precis de som säkerhetskravet skrevs för. Vad säger det dig om
    värdet av en sanningstabell du härlett ur ett krav, till skillnad från en du läst av ur en
    krets du redan byggt?
\end{itemize}

\modulefig[0.72]{lectures/L01/appendix/images/adas.png}

\begin{selfcheck}
\textbf{Självkontroll.} Döp din entitet till \code{adas}, med ingångarna \code{driver_brake},
\code{sensor}, \code{radar}, \code{error} och utgången \code{engine_brake}, deklarerade i den
ordningen; dess testbänk ligger i \file{exercises/adas/} och kontrollerar alla sexton raderna i
\secref{c1:sec:adas}:s sanningstabell.
\end{selfcheck}
